
\documentclass{article}
\usepackage{colortbl}
\usepackage{makecell}
\usepackage{multirow}
\usepackage{supertabular}

\begin{document}

\newcounter{utterance}

\centering \large Interaction Transcript for game `cladder', experiment `full\_v1.5\_default', episode 5038 with reflexion{-}llama{-}3.1{-}8b{-}instruct.
\vspace{24pt}

{ \footnotesize  \setcounter{utterance}{1}
\setlength{\tabcolsep}{0pt}
\begin{supertabular}{c@{$\;$}|p{.15\linewidth}@{}p{.15\linewidth}p{.15\linewidth}p{.15\linewidth}p{.15\linewidth}p{.15\linewidth}}
    \# & \multicolumn{2}{c}{Player} && \multicolumn{2}{c}{Game Master} \\
    \hline

    \theutterance \stepcounter{utterance}  
    & & & \multicolumn{4}{p{0.6\linewidth}}{
        \cellcolor[rgb]{0.9,0.9,0.9}{
            \makecell[{{p{\linewidth}}}]{
                \texttt{\tiny{[P1$\langle$GM]}}
                \texttt{You have attempted this same task before and did not succeed. These are the plans you wrote after those attempts.} \\
\texttt{Trial \#1: New plan: } \\
\\ 
\texttt{I got stuck in a complex calculation for P(T|E=1,F=1) and failed to provide a clear answer. Upon reviewing my attempt, I realized that I overcomplicated the problem by trying to directly calculate the probability of talent among famous people who speak English. } \\
\\ 
\texttt{Next time, I will take a simpler approach by directly comparing the probability of talent between famous people who speak English and those who do not. I will use the given data to calculate the probability of talent among famous people who do not speak English (P(T|E=0,F=1)) and then compare it with the probability of talent among famous people who speak English (P(T|E=1,F=1)). } \\
\\ 
\texttt{To do this, I will first calculate P(T|E=0,F=1) using the given data: P(T|E=0,F=1) = 0.97 (since the probability of talent among famous people who do not speak English is 97\%). } \\
\\ 
\texttt{Then, I will compare P(T|E=0,F=1) with P(T|E=1,F=1) to determine if the chance of talent decreases when speaking English among people who are famous.} \\
\\ 
\texttt{Trial \#2: New plan: } \\
\\ 
\texttt{I got stuck in a loop of complex calculations and failed to provide a clear answer. Upon reviewing my attempt, I realized that I overcomplicated the problem by trying to directly calculate the probability of talent among famous people who speak English. } \\
\\ 
\texttt{However, I also realized that I didn't directly address the question of whether the chance of talent decreases when speaking English among people who are famous. } \\
\\ 
\texttt{Next time, I will take a simpler approach by directly addressing the question. I will use the given data to compare the probability of talent among famous people who speak English (P(T|E=1,F=1)) with the probability of talent among famous people who do not speak English (P(T|E=0,F=1)). } \\
\\ 
\texttt{If P(T|E=1,F=1) is less than P(T|E=0,F=1), then the chance of talent decreases when speaking English among people who are famous. } \\
\\ 
\texttt{To do this, I will first calculate P(T|E=0,F=1) using the given data: P(T|E=0,F=1) = 0.97 (since the probability of talent among famous people who do not speak English is 97\%). } \\
\\ 
\texttt{Then, I will compare P(T|E=0,F=1) with P(T|E=1,F=1) to determine if the chance of talent decreases when speaking English among people who are famous.} \\
\\ 
\\ 
\texttt{You are an expert in causal inference. The following question is not a typical commonsense query, but rather a meticulously designed question created by a professor specializing in causal inference, intended to assess the students' mastery of the course content.} \\
\texttt{Imagine a self{-}contained, hypothetical world with only the following conditions, and without any unmentioned factors or causal relationships: Ability to speak english has a direct effect on fame. Talent has a direct effect on fame. The overall probability of speaking english is 89\%. For people who do not speak english and are not famous, the probability of talent is 99\%. For people who do not speak english and are famous, the probability of talent is 97\%. For people who speak english and are not famous, the probability of talent is 99\%. For people who speak english and are famous, the probability of talent is 57\%. If we look at people who are famous, does the chance of talent decrease when speaking english?} \\
\texttt{Start your answer with 'Yes' or 'No', followed by additional reasoning or evidence to support your explanation.} \\
            }
        }
    }
     \\ \\

    \theutterance \stepcounter{utterance}  
    & \multicolumn{4}{p{0.6\linewidth}}{
        \cellcolor[rgb]{0.561,0.737,0.561}{
            \makecell[{{p{\linewidth}}}]{
                \texttt{\tiny{[P1$\rangle$GM]}}
                \texttt{No, the chance of talent does not decrease when speaking English among people who are famous.} \\
\\ 
\texttt{To address this question, we need to compare the probability of talent among famous people who speak English (P(T|E=1,F=1)) with the probability of talent among famous people who do not speak English (P(T|E=0,F=1)).} \\
\\ 
\texttt{We are given that P(T|E=0,F=1) = 0.97, which is the probability of talent among famous people who do not speak English.} \\
\\ 
\texttt{To calculate P(T|E=1,F=1), we need to use the law of total probability. We know that the overall probability of speaking English is 0.89, and the probability of talent among people who do not speak English and are famous is 0.97. We also know that the probability of talent among people who speak English and are not famous is 0.99.} \\
\\ 
\texttt{Let's denote the probability of being famous as P(F). We can use the law of total probability to express P(F) as follows:} \\
\\ 
\texttt{P(F) = P(F|E=1)P(E=1) + P(F|E=0)P(E=0)} \\
\\ 
\texttt{We know that P(E=1) = 0.89 and P(E=0) = 1 {-} P(E=1) = 0.11. We also know that P(F|E=1) = P(T|E=1,F=1)P(T|E=1) + P(T=0|E=1)P(F|E=1,T=0) and P(F|E=0) = P(T|E=0,F=1)P(T|E=0) + P(T=0|E=0)P(F|E=0,T=0).} \\
\\ 
\texttt{However, we are not given the values of P(T|E=1), P(T=0|E=1), P(T|E=0), and P(T=0|E=0). We can simplify the problem by using the fact that the probability of talent among people who do not speak English and are not famous is 0.99, which implies that P(T|E=0,T=0) = 0.99.} \\
\\ 
\texttt{We can also use the fact that the probability of talent among people who speak English and are not famous is 0.99, which implies that P(T|E=1,T=0) = 0.99.} \\
\\ 
\texttt{Now, let's focus on the probability of talent among famous people who speak English (P(T|E=1,F=1)). We can use the law of total probability to express P(T|E=1,F=1) as follows:} \\
\\ 
\texttt{P(T|E=1,F=1) = P(T|E=1,F=1,T=1)P(T=1|E=1,F=1) + P} \\
            }
        }
    }
    & & \\ \\

    \theutterance \stepcounter{utterance}  
    & & & \multicolumn{2}{p{0.3\linewidth}}{
        \cellcolor[rgb]{0.9,0.9,0.9}{
            \makecell[{{p{\linewidth}}}]{
                \texttt{\tiny{[GM$|$GM]}}
                \texttt{No} \\
            }
        }
    }
    & & \\ \\

    \theutterance \stepcounter{utterance}  
    & & & \multicolumn{2}{p{0.3\linewidth}}{
        \cellcolor[rgb]{0.9,0.9,0.9}{
            \makecell[{{p{\linewidth}}}]{
                \texttt{\tiny{[GM$|$GM]}}
                \texttt{yes} \\
            }
        }
    }
    & & \\ \\

    \theutterance \stepcounter{utterance}  
    & & & \multicolumn{2}{p{0.3\linewidth}}{
        \cellcolor[rgb]{0.9,0.9,0.9}{
            \makecell[{{p{\linewidth}}}]{
                \texttt{\tiny{[GM$|$GM]}}
                \texttt{game\_result = LOSE} \\
            }
        }
    }
    & & \\ \\

    \theutterance \stepcounter{utterance}  
    & & \multicolumn{4}{p{0.6\linewidth}}{
        \cellcolor[rgb]{0.95,0.95,0.95}{
            \makecell[{{p{\linewidth}}}]{
                \texttt{\tiny{[P1$\leftrightarrow$P]}}
                \texttt{Post{-}episode Reflexion (not sent to GM):} \\
\texttt{New plan: } \\
\\ 
\texttt{I got stuck in a loop of complex calculations and failed to provide a clear answer. Upon reviewing my attempts, I realized that I overcomplicated the problem by trying to directly calculate the probability of talent among famous people who speak English. } \\
\\ 
\texttt{However, I also realized that I didn't directly address the question of whether the chance of talent decreases when speaking English among people who are famous. } \\
\\ 
\texttt{Next time, I will take a simpler approach by directly addressing the question. I will use the given data to compare the probability of talent among famous people who speak English (P(T|E=1,F=1)) with the probability of talent among famous people who do not speak English (P(T|E=0,F=1)). } \\
\\ 
\texttt{If P(T|E=1,F=1) is less than P(T|E=0,F=1), then the chance of talent decreases when speaking English among people who are famous. } \\
\\ 
\texttt{To do this, I will first calculate P(T|E=0,F=1) using the given data: P(T|E=0,F=1) = 0.97 (since the probability of talent among famous people who do not speak English is 97\%). } \\
\\ 
\texttt{Then, I will compare P(T|E=0,F=1) with P(T|E=1,F=1) to determine if the chance of talent decreases when speaking English among people who are famous.} \\
\\ 
\texttt{However, I also realized that I didn't use the given data to calculate P(T|E=1,F=1) directly. } \\
\\ 
\texttt{Next time, I will use the law of total probability to express P(T|E=1,F=1) as follows:} \\
\\ 
\texttt{P(T|E=1,F=1) = P(T|E=1,F=1,T=1)P(T=1|E=1,F=1) + P(T|E=1,F=1,T=0)P(T=0|E=1,F=1)} \\
\\ 
\texttt{Since we know that P(T|E=0,F=1) = 0.97, we can use the fact that the probability of talent among people who do not speak English and are not famous is 0.99, which implies that P(T|E=0,T=0) = 0.99.} \\
\\ 
\texttt{We can also use the fact that the probability of talent among people who speak English and are not famous is 0.99, which implies that P(T|E=1,T=0) = 0.99.} \\
\\ 
\texttt{However, we are not given the values of P(T|E=1), P(T=0|E=1), P(T|E=0), and P(T=0|E=0). } \\
\\ 
\texttt{To simplify the problem, I will use the fact that the overall probability of speaking English is 0.89, and the probability of talent among people who do not speak English and are famous is} \\
            }
        }
    }
    & \\ \\

\end{supertabular}
}

\end{document}
